% 样本需求与标注者能力的候选关系、实际分配和弃权转交。
\begin{tikzpicture}[x=1mm,y=1mm,font=\figurefont\footnotesize,text=BookInk,
  sample/.style={draw=FigureInput,fill=FigureInput!10,line width=.9pt,rounded corners=1.2mm,minimum width=38mm,minimum height=12mm,align=center,inner sep=2.5pt},
  worker/.style={draw=FigureModel,fill=FigureModel!10,line width=.9pt,rounded corners=1.2mm,minimum width=39mm,minimum height=12mm,align=center,inner sep=2.5pt},
  candidate/.style={draw=BookMuted!65,line width=.65pt},
  assigned/.style={draw=BookTeal,line width=1.5pt},
  transfer/.style={knowledge arrow,draw=BookAmber,line width=.9pt,dashed,rounded corners=1.2mm}]
  \node[font=\figurefont\bfseries\small] at (20,51) {样本需求};
  \node[font=\figurefont\bfseries\small] at (89,51) {标注者能力与负载};
  \node[sample] (s1) at (20,36) {普通中文文档\\低难度};
  \node[sample] (s2) at (20,16) {医疗证据片段\\需查阅术语};
  \node[sample] (s3) at (20,-4) {英文合同条款\\高难度};
  \node[worker] (w1) at (89,36) {受训通用标注者\\中文，余6项};
  \node[worker] (w2) at (89,16) {医疗领域专家\\中文，余1项};
  \node[worker] (w3) at (89,-4) {双语法律人员\\中英，余2项};
  \draw[candidate] (s1.east)--(w1.west);
  \draw[candidate] (s2.east)--(w1.west);
  \draw[candidate] (s2.east)--(w2.west);
  \draw[candidate] (s3.east)--(w3.west);
  \draw[assigned] (s1.east)--(w1.west);
  \draw[assigned] (s2.east)--(w2.west);
  \draw[assigned] (s3.east)--(w3.west);
  \draw[transfer] (w1.east)--(112,36)--node[left,align=right] {难例弃权\\转交}(112,16)--(w2.east);
  \draw[candidate] (8,-18)--(27,-18);
  \node[anchor=west] at (29,-18) {可承接};
  \draw[assigned] (58,-18)--(77,-18);
  \node[anchor=west] at (79,-18) {本轮分配};
  \draw[transfer] (8,-27)--(27,-27);
  \node[anchor=west] at (29,-27) {来源弃权后的转交};
\end{tikzpicture}
